% Part: computability
% Chapter: computability-theory
% Section: non-comp-set

\documentclass[../../../include/open-logic-section]{subfiles}

\begin{document}

\olfileid[hi]{cmp}{thy}{ncp}
\olsection{असंगणनीय समुच्चय अस्तित्व में हैं}


हमने ऊपर देखा कि प्रत्येक संगणनीय समुच्चय संगणनीयतः परिगणनीय है। क्या
इसका विलोम सत्य है? अगला परिणाम दिखाता है कि सामान्यतः ऐसा नहीं है।

\begin{thm}\ollabel{thm:K-0}
$K_0$ को समुच्चय $\Setabs{\tuple{e, x}}{\cfind{e}(x) \fdefined}$ मानें।
तब $K_0$ संगणनीयतः परिगणनीय है, किंतु संगणनीय नहीं।
\end{thm}

\begin{proof}
$K_0$ को संगणनीयतः परिगणनीय देखने के लिए ध्यान दें कि वह इस प्रकार
परिभाषित फलन~$f$ का प्रांत है:
\[
f(z) = \umin{y}{(\len{z} = 2 \land T((z)_0, (z)_1, y))}.
\]
वास्तव में, यदि $\cfind{e}(x)$ परिभाषित है, तो $f(\tuple{e, x})$ रुकने
वाला कोई संगणना अनुक्रम खोज लेता है; यदि $\cfind{e}(x)$ अपरिभाषित है,
तो $f(\tuple{e, x})$ भी अपरिभाषित है; और यदि $z$ किसी युग्म को कूटित तक
नहीं करता, तो $f(z)$ भी अपरिभाषित है।

$K_0$ का असंगणनीय होना ठीक विराम समस्या की अनिर्णेयता,
\olref[hlt]{thm:halting-problem}, है।
\end{proof}

समुच्चय $K_0$ उन युग्मों $\tuple{e,x}$ का समुच्चय है जिनके लिए
$\cfind{e}(x) \fdefined$; अर्थात् $\tuple{e,x} \in K_0$ तभी और केवल
तभी जब $\cfind{e}$, आगत~$x$ पर परिभाषित है (रुकता है), इसलिए इसे
``विराम समुच्चय'' भी कहते हैं। समुच्चय
$K = \Setabs{e}{\cfind{e}(e) \fdefined}$ ``स्व-विराम समुच्चय'' है। इसे
प्रायः एक आदर्श अनिर्णेय समुच्चय के रूप में प्रयोग किया जाता है।

\begin{thm}\ollabel{thm:K}
स्व-विराम समुच्चय $K = \Setabs{e}{\cfind{e}(e) \fdefined}$
!!{c.e.} है, किंतु निर्णेय नहीं।
\end{thm}

\begin{proof}
  मान लें कि $K$ निर्णेय है, अर्थात् उसका अभिलाक्षणिक फलन
  $\Char{K}$ संगणनीय है। मानें
  \[d(e) = \begin{cases}
  1 & \text{यदि\/ $\Char{K}(e) = 0$}\\
  \fundefined & \text{अन्यथा।}
  \end{cases}
  \] 
  $k$ को~$d$ का सूचकांक मानें, अर्थात् $d \simeq \cfind{k}$। तब $d(k)
  \simeq \cfind{k}(k)$। यह इस तथ्य के विरुद्ध है कि $d(k)
  \fdefined$
  तभी और केवल तभी जब $\cfind{k}(k) \fundefined$; यह~$d$ की परिभाषा से
  निष्पन्न होता है।

  $K$, $f(x) = \umin{y}{T(x,x,y)}$ का प्रांत है और इसलिए !!{c.e.} है।
\end{proof}

\end{document}
